% Quarantined v3 strict-open repair candidate for P06.2.
% This file retains the v2 arguments that passed production review and deletes
% the two endpoint extensions rejected there.  It is not authorized for live
% integration before a new exact-byte hostile audit and a new production
% Rethlas verification.

\begin{proof}[Proof of \cref{prop:p062-product-fibre-zero-current-candidate}]
Write \(m=\dim B\) and let
\(
 \pi\colon X=B\times_KZ\to B
\)
be the projection. It is flat, boundaryless, and has pure
one-dimensional fibres. Fix \(x\in B\) and a compactly supported ordinary
CLD-smooth function \(\varphi\) on \(X_x=Z_{\mathscr H(x)}\). We first
record why, after shrinking \(B\) around \(x\), this test function may be
spread to the product.

Choose an affinoid neighbourhood \(U_0\) of \(x\) in \(B\), and put
\[
                 Y_0=U_0\times_K Z,
                 \qquad F=(Y_0)_x=X_x.
\]
The space \(Y_0\) is good and topologically separated. In particular it
is almost good by \cite[\S11.1.10]{CLD25}, and it is locally
holomorphically separated by \cite[\S\S15.2.2--15.2.4]{CLD25}.

Fix \(y\in\operatorname{supp}(\varphi)\). By the local moment
description of ordinary smooth functions
\cite[Corollary~15.1.2]{CLD25}, there are a neighbourhood \(V_y'\) of
\(y\) in \(F\), invertible \(\mathscr H(x)\)-analytic functions
\(f_{y,1},\ldots,f_{y,N_y}\) on \(V_y'\), and a smooth real function
\(\rho_y\), defined on a neighbourhood of their tropical image, such
that
\[
 \varphi=\rho_y\bigl(-\log|f_{y,1}|,\ldots,
                         -\log|f_{y,N_y}|\bigr)
 \qquad\text{on }V_y'.
\]
By the analytic-fibre construction of \cite[\S11.1.10]{CLD25}, the
completed residue fields at \(y\) on \(F\) and on \(Y_0\) are canonically
identified; write their common value as \(\mathscr H_{Y_0}(y)\). For
completeness, consider affinoid neighbourhoods associated with a map
\(A\to C\), the fibre point is induced on
\(
 \mathcal M(C\widehat\otimes_A\mathscr H(x))
\)
by evaluation in \(\mathscr H_{Y_0}(y)\). The fraction field of the image
of \(C\) is already dense in \(\mathscr H_{Y_0}(y)\), while the image of
\(\mathscr H(x)\) lies in its closure. Hence the completed residue field
of the induced fibre point is exactly \(\mathscr H_{Y_0}(y)\).

Since \(Y_0\) is good at \(y\), \cite[\S11.1.10]{CLD25} gives a dense
embedding
\[
                  \kappa_{Y_0}(y)\lhook\joinrel\longrightarrow
                  \mathscr H_{Y_0}(y).
\]
For every \(j\), choose \(c_{y,j}\in\kappa_{Y_0}(y)\) such that
\begin{equation}\label{eq:p062-v2-strict-moment-approximation}
 |c_{y,j}-f_{y,j}(y)|<|f_{y,j}(y)|.
\end{equation}
Represent \(c_{y,j}\) by an analytic germ
\(g_{y,j}\in\mathcal O_{Y_0,y}\). After restriction to \(F\), continuity
of absolute values preserves the strict inequality in
\eqref{eq:p062-v2-strict-moment-approximation} on a sufficiently small
neighbourhood \(V_y\subseteq V_y'\). The ultrametric inequality therefore
gives the exact identities
\begin{equation}\label{eq:p062-v2-exact-moment-norms}
             |g_{y,j}|_{F}=|f_{y,j}|\qquad(1\leq j\leq N_y)
             \quad\text{on }V_y.
\end{equation}
The germs \(g_{y,j}\) are units at \(y\). Shrinking an ambient ordinary
neighbourhood \(W_y\) of \(y\), we may assume that they are represented
by invertible analytic functions on \(W_y\), that
\(W_y\cap F\subseteq V_y\), and that their tropical vector lies in the
domain of \(\rho_y\). Thus
\[
 \Phi_y=\rho_y\bigl(-\log|g_{y,1}|,\ldots,
                         -\log|g_{y,N_y}|\bigr)
\]
is ordinary CLD-smooth on \(W_y\), and
\begin{equation}\label{eq:p062-v2-local-fibre-extension}
                      \Phi_y|_{W_y\cap F}=\varphi|_{W_y\cap F}
\end{equation}
literally.

Put \(S=\operatorname{supp}(\varphi)\). It is a compact closed subset of
the Hausdorff space \(Y_0\). Choose finitely many of the preceding
neighbourhoods, say \(W_1,\ldots,W_N\), which cover \(S\). By
\cite[Corollary~15.2.9]{CLD25}, there are nonnegative ordinary smooth
functions \(\theta_1,\ldots,\theta_N\) on \(Y_0\) such that
\[
 \begin{gathered}
  \operatorname{supp}(\theta_i)\text{ is compact and contained in }W_i,\\
  \sum_{i=1}^N\theta_i=1\quad\text{on a neighbourhood of }S.
 \end{gathered}
\]
The product \(\theta_i\Phi_i\), formed on \(W_i\), extends by zero to an
ordinary smooth function on \(Y_0\). Define
\[
                         \Phi=\sum_{i=1}^N\theta_i\Phi_i.
\]
If \(z\in F\), every summand which can be nonzero satisfies
\(\Phi_i(z)=\varphi(z)\) by
\eqref{eq:p062-v2-local-fibre-extension}. Hence
\[
                         \Phi(z)=\varphi(z)
                                      \sum_{i=1}^N\theta_i(z).
\]
For \(z\in S\) the sum is one, while for \(z\notin S\) one has
\(\varphi(z)=0\). Therefore \(\Phi|_F=\varphi\) everywhere. Moreover,
\(
 \operatorname{supp}(\Phi)\subseteq
 \bigcup_i\operatorname{supp}(\theta_i)
\)
is compact, so its projection to the Hausdorff affinoid \(U_0\) is
proper. Taking \(U=U_0\), we have constructed an ordinary CLD-smooth
function \(\Phi\) on \(U\times_KZ\), with support proper over \(U\), such
that
\begin{equation}\label{eq:p062-spread-test-candidate}
                         \Phi|_{X_x}=\varphi.
\end{equation}
This construction uses ambient smooth germs and an ambient partition; it
does not spread rational domains and does not use a pullback of a current.

After replacing \(U\) by a relatively compact affinoid neighbourhood of
\(x\), put
\[
 \tau=\mathrm d'\mathrm d''\Phi,
 \qquad
 G(t)=\int_{X_t}h_t\,\tau|_{X_t},
 \qquad t\in U.
\]
Although an arbitrary continuous \(h\) need not be a G-locally tropical
coefficient, \(G\) is well defined and continuous. Here are the details
which avoid that coefficient error. On a compact neighbourhood of
\(\operatorname{supp}(\tau)\), choose ordinary CLD-smooth functions
\(h_\nu\) converging uniformly to \(h\)
\cite[Theorem~15.2.6]{CLD25}, and set
\[
                         G_\nu(t)=
                 \int_{X_t}h_\nu\,\tau|_{X_t}.
\]
The form \(h_\nu\tau\) is smooth of relative top degree and has
\(\pi\)-compact support. Fibre integration therefore makes \(G_\nu\)
G-smooth, hence continuous. By the decomposition of a compactly
supported symmetric smooth form as a difference of two compactly
supported positive forms \cite[Proposition~6.5.2(b)]{CLD25}, write
\(
 \tau=\tau^+-\tau^-.
\)
Fibre integration makes the two functions
\(
 t\mapsto\int_{X_t}\tau^\pm|_{X_t}
\)
continuous by \cite[Theorem~16.3.3(d)]{CLD25}, hence bounded on the
compact base under consideration.
Consequently
\begin{equation}\label{eq:p062-uniform-fibre-variation-candidate}
 \sup_{t\in U}|G_\nu(t)-G(t)|
 \leq \lVert h_\nu-h\rVert_\infty
       \sup_{t\in U}\left(
       \int_{X_t}\tau^+ +\int_{X_t}\tau^-\right)\longrightarrow0.
\end{equation}
Thus \(G\) is continuous. The same approximation, the projection formula
\cite[Theorem~16.5.4 and Corollary~16.5.5]{CLD25}
for each \(h_\nu\tau\), and
\eqref{eq:p062-uniform-fibre-variation-candidate} give, for every compactly
supported smooth top-degree form \(\alpha\) on \(U\),
\begin{equation}\label{eq:p062-continuous-fubini-candidate}
 \int_{U\times_KZ}h\,\tau\wedge\pi^*\alpha
                         =\int_U G\,\alpha.
\end{equation}

The form \(\Phi\pi^*\alpha\) has bidegree
\((m,m)=(\dim X-1,\dim X-1)\) and compact support. Since \(\alpha\) has
top degree on the \(m\)-dimensional base,
\(
 \mathrm d'\alpha=\mathrm d''\alpha=0
\), and hence
\[
 \mathrm d'\mathrm d''(\Phi\pi^*\alpha)
        =(\mathrm d'\mathrm d''\Phi)\wedge\pi^*\alpha.
\]
Using \eqref{eq:p062-product-zero-current-candidate} and
\eqref{eq:p062-continuous-fubini-candidate}, we obtain
\[
                              \int_U G\,\alpha=0
\]
for every such \(\alpha\). A standard strict polycylinder is very rich
\cite[Theorem~16.4.4(b)]{CLD25},
so the map which sends a continuous function to its associated
\((0,0)\)-current is injective
\cite[Corollary~7.3.12]{CLD25}. Therefore \(G=0\) on \(U\). Evaluating at
\(x\) and using \eqref{eq:p062-spread-test-candidate} gives
\[
 0=G(x)=\int_{X_x}h_x\,\mathrm d'\mathrm d''\varphi
       =\bigl\langle\mathrm d'\mathrm d''[h_x],\varphi\bigr\rangle.
\]
Since \(\varphi\) was arbitrary, this proves
\eqref{eq:p062-fibre-zero-current-candidate}.
\end{proof}

\begin{proof}[Proof of \cref{prop:p062-arbitrary-field-annulus-candidate}]
For \(s\in(r,R)\), let \(\eta_s\) be the \(s\)-Gauss point. Testing the
current equation on compactly supported functions of
\(-\log|T|\), and using the weight-one split-torus integration formula,
shows that
\begin{equation}\label{eq:p062-annulus-radial-affine-candidate}
             u\longmapsto h(\eta_{\mathrm e^{-u}})
\end{equation}
has zero distributional second derivative on
\((-\log R,-\log r)\). Thus it is affine. This is exactly the
one-dimensional part of the split-torus calculation used in
\cref{prop:p062-rational-centred-traces}; unlike the disc case, there is
no endpoint at \(T=0\) and its affine slope is not forced to vanish.

We next prove that every rigid point has the value prescribed by
\eqref{eq:p062-annulus-radial-affine-candidate}. Let \(y\) be rigid, put
\(E=\mathscr H(y)\), \(s=|T(y)|\), and let
\[
                 P(T)=T^d+\sum_{i=0}^{d-1}c_iT^i
\]
be the monic irreducible polynomial of \(T(y)\). Completeness of \(K\)
gives a unique extension of its absolute value to every finite extension.
Every conjugate of \(T(y)\), counted with inseparable multiplicity, has
absolute value \(s\). Consequently
\begin{equation}\label{eq:p062-annulus-coefficient-bounds-candidate}
             |c_i|\leq s^{d-i},
             \qquad |c_0|=s^d=:S\in|K^\times|.
\end{equation}
Strict dominance of the constant term for \(|T|<s\) and of the leading
term for \(|T|>s\) gives
\begin{equation}\label{eq:p062-annulus-polynomial-tube-candidate}
 X_y:=\{|P(T)|<S\}\subseteq\{|T|=s\}\subset A^-_{r,R}.
\end{equation}
The polynomial \(P\) induces a finite flat morphism of degree \(d\)
\[
             p\colon X_y\longrightarrow D^-_S,
             \qquad Z=P(T).
\]
Let \(H=\operatorname{Tr}_p(h|_{X_y})\), with analytic local degrees in
the fibre trace. By \cref{prop:p062-finite-flat-continuous-trace}, \(H\)
is continuous and \(\mathrm d'\mathrm d''[H]=0\). The arbitrary-field
strict-disc theorem \cref{prop:p062-arbitrary-field-disc} makes \(H\)
constant. The zero fibre is the single analytic point \(y\), with degree
\(d\), so
\begin{equation}\label{eq:p062-annulus-trace-zero-candidate}
                              H=d\,h(y).
\end{equation}

Let \(\xi_S\) denote the boundary Gauss point of the target closed disc.
The graded Gauss calculation used in the proof of
\cref{prop:p062-arbitrary-field-disc}, now at \(t=s\), gives
\begin{equation}\label{eq:p062-annulus-single-gauss-fibre-candidate}
                      P^{-1}(\xi_S)=\{\eta_s\},
                      \qquad \deg_{\eta_s}(P)=d.
\end{equation}
Indeed, the initial polynomial
\(\operatorname{in}_s(P)(U)\) has degree \(d\); the rational-function
extension generated by \(U\) over that generated by
\(\operatorname{in}_s(P)(U)\) has degree \(d\), including its inseparable
part. The graded fundamental inequality gives local degree at least \(d\),
and finite flatness leaves neither degree nor another point in the fibre.

Choose a neighbourhood \(W\) of \(\eta_s\) whose closure is contained in
the annulus. Since \(P\) is finite and
\eqref{eq:p062-annulus-single-gauss-fibre-candidate} is a singleton fibre,
there is an analytic neighbourhood \(V\) of \(\xi_S\) such that
\(P^{-1}(V)\subset W\). Applying the continuity part of
\cref{prop:p062-finite-flat-continuous-trace} to
\(P^{-1}(V)\to V\) shows, as \(q\uparrow S\), that
\[
             H(\xi_q)\longrightarrow d\,h(\eta_s).
\]
Together with \eqref{eq:p062-annulus-trace-zero-candidate}, this yields
\begin{equation}\label{eq:p062-annulus-rigid-to-skeleton-candidate}
                              h(y)=h(\eta_s).
\end{equation}
Rigid points are dense locally on the strict annulus. The affinity of
\eqref{eq:p062-annulus-radial-affine-candidate},
\eqref{eq:p062-annulus-rigid-to-skeleton-candidate}, and continuity give
\eqref{eq:p062-annulus-affine-candidate} at every analytic point.
Uniqueness is immediate because \(\log|T|\) is nonconstant on the
skeleton.
\end{proof}

\begin{proof}[Proof of \cref{thm:p062-arbitrary-field-polycylinder-candidate}]
We first work with the open polycylinder. Apply
\cref{prop:p062-product-fibre-zero-current-candidate} to any coordinate
projection. On a disc fibre, \cref{prop:p062-arbitrary-field-disc} shows
that the restriction of \(h\) is constant. Hence \(h\) is independent of
every disc coordinate.
If \(q=0\), this already proves that \(h\) is constant, and there is
nothing more to show.

Fixing the \(K\)-rational origin in every disc factor and replacing the
disc coordinates successively shows, by the preceding constancy, that
\(h\) is the pullback of a continuous function on the product of the
annulus factors. We continue to denote that function by \(h\).

On an annulus fibre, \cref{prop:p062-arbitrary-field-annulus-candidate}
shows that the restriction of \(h\) is affine in the logarithmic radius.
In particular, it has at an arbitrary fibre point the same value as at the
Gauss point with the same radius. The slicing proposition is valid over
every base point, with its possibly nonrigid complete residue field. It
remains to verify that successive replacements in distinct coordinates
have a common, order-independent endpoint.

Fix radii \(\mathbf s=(s_1,\ldots,s_q)\) in the open radius rectangle. For
a Laurent polynomial
\(
 f=\sum_{\nu\in\mathbf Z^q}a_\nu S^\nu
\), successive one-variable Gauss extensions, in any order, give
\begin{equation}\label{eq:p062-iterated-gauss-v2}
                  |f|_{\mathbf s}
                   =\max_{\nu}|a_\nu|
                         \prod_{j=1}^q s_j^{\nu_j}.
\end{equation}
Indeed, each step takes the maximum over one exponent and the finitely
many maximum operations commute. At an intermediate step the coefficients
may lie in the completed residue field of the coordinates not yet
replaced. The one-variable Gauss norm is still defined by the same maximum
formula, and passing to that completion does not change the result.

The formula extends from Laurent polynomials to every Laurent series
converging on a fixed compact subannulus. Such a series is the Banach-norm
limit of finite Laurent sums, and every bounded multiplicative seminorm is
continuous for that norm. It follows that every order of replacement
induces the same seminorm on each compact subannulus, hence the same
analytic point on the open polyannulus. By
\eqref{eq:p062-iterated-gauss-v2}, that point is the canonical
multivariate Gauss point \(\sigma(\mathbf s)\).

The section \(\mathbf s\mapsto\sigma(\mathbf s)\) is continuous. Indeed,
on a compact logarithmic subrectangle, the norm of a Laurent polynomial is
the maximum of finitely many continuous monomials in \(\mathbf s\), and
the preceding Banach approximation gives the norm of every analytic
function as a locally uniform limit of such functions. This is precisely
continuity for the Berkovich topology.

Now let \(z\) be an arbitrary point of the open polyannulus and put
\(s_j=|S_j(z)|\). Replacing the coordinates one at a time does not change
the value of \(h\), while the preceding calculation identifies the final
point independently of order. Therefore
\[
                     h(z)=h\bigl(\sigma(\mathbf s)\bigr).
\]
It follows that \(h\) factors through the standard tropicalization:
\begin{equation}\label{eq:p062-polycylinder-factor-candidate}
 h=\phi\circ F,
 \qquad
 F=(-\log|S_1|,\ldots,-\log|S_q|),
\end{equation}
where
\[
             \phi(\mathbf u)
                =h\bigl(\sigma(\mathrm e^{-u_1},\ldots,
                                      \mathrm e^{-u_q})\bigr).
\]
The continuity just proved for the Gauss section and continuity of \(h\)
show that \(\phi\) is continuous on the open rectangle
\(
 \prod_j(-\log R'_j,-\log r_j).
\)

If \(p>0\), choose a compactly supported smooth positive top-degree form
\(\gamma\) in the product of the disc factors with
\(\int\gamma>0\). If \(p=0\), put \(\gamma=1\) and
\(\int\gamma=1\). For an arbitrary compactly supported Lagerberg form
\(\beta\) of bidegree \((q-1,q-1)\) on the logarithmic rectangle, use as
a test form the product of \(\gamma\) with the pullback of \(\beta\).
Since \(\gamma\) has top degree in the disc variables,
\(\mathrm d'\gamma=\mathrm d''\gamma=0\). The weight-one split-torus
integration formula and Fubini therefore convert
\(
 \mathrm d'\mathrm d''[h]=0
\)
into
\(
 \mathrm d'\mathrm d''[\phi]=0
\)
in the distributional Lagerberg sense on that rectangle; the nonzero
factor \(\int\gamma\) cancels. Equivalently,
every distribution
\(
 \partial^2\phi/\partial u_i\partial u_j
\)
vanishes. Convolution on relatively compact Euclidean balls makes every
regularization affine; compatibility on overlaps and continuity show that
\(\phi\) itself is affine. This proves
\eqref{eq:p062-polycylinder-affine-candidate}.

Each \(\log|S_j|\) is the logarithm of the norm of an analytic unit.
Thus the displayed function is ordinary CLD-smooth with zero Hessian,
weakly smooth and Gubler--Jell--Rabinoff harmonic, and classically
pluriharmonic. Uniqueness follows by restriction to the standard
polyannulus skeleton.
\end{proof}
